\documentclass[11pt,a4paper]{article}
\usepackage[T1]{fontenc}
\usepackage[utf8]{inputenc}
\usepackage{lmodern,microtype}
\usepackage[margin=.85in]{geometry}
\usepackage{amsmath,amssymb,amsthm,mathtools,aliascnt}
\usepackage{booktabs,tabularx,array,enumitem,needspace,float}
\usepackage[section]{placeins}
\usepackage{xcolor,tikz,mdframed}
\usetikzlibrary{arrows.meta,positioning,calc}
\usepackage[hidelinks]{hyperref}
\usepackage[nameinlink,noabbrev]{cleveref}
\usepackage{fancyhdr}
\hypersetup{pdftitle={Sync HotStuff RR v4.2},pdfsubject={Regional finality and L1-authorized external truncation}}
\setlength{\headheight}{14pt}
\pagestyle{fancy}\fancyhf{}
\lhead{\small Sync HotStuff RR v4.2}\rhead{\small L1-authorized regional sessions}
\cfoot{\thepage}
\setlength{\parindent}{0pt}\setlength{\parskip}{4pt}
\setlength{\emergencystretch}{2em}
\setlist{nosep,leftmargin=*}
\newcolumntype{Y}{>{\raggedright\arraybackslash}X}
\definecolor{rrRed}{HTML}{A51C30}
\definecolor{rrBlue}{HTML}{2563EB}
\definecolor{rrTeal}{HTML}{0F766E}
\newcommand{\RR}[1]{{\color{rrRed}#1}}
\newcommand{\QC}{\mathcal C}
\newcommand{\FC}{\mathcal F}
\newcommand{\Hash}{\mathsf H}
\newcommand{\rank}{\operatorname{rank}}
\newcommand{\wt}{\operatorname{wt}}
\newcommand{\tagg}[1]{\ensuremath{\mathsf{#1}}}
\newtheorem{definition}{Definition}[section]
\newaliascnt{assumption}{definition}
\newtheorem{assumption}[assumption]{Assumption}
\aliascntresetthe{assumption}
\crefname{assumption}{assumption}{assumptions}
\Crefname{assumption}{Assumption}{Assumptions}
\newcounter{protocol}
\newenvironment{protocol}[2]{\refstepcounter{protocol}\label{#1}
\begin{mdframed}[nobreak=true,linewidth=.6pt,skipabove=8pt,skipbelow=8pt,
innertopmargin=7pt,innerbottommargin=7pt]\small
\textbf{Protocol \theprotocol. #2}\par\smallskip}{\end{mdframed}}
\crefname{protocol}{Protocol}{Protocols}
\Crefname{protocol}{Protocol}{Protocols}
\title{Sync HotStuff RR v4.2\\\large Regional Finality and L1-Authorized Rotation}
\author{}\date{12 October 2026}
\begin{document}
\maketitle
\begin{abstract}
L1 supplies stake configurations, leader schedules, regional quotas, and
status updates. Within each region, Sync HotStuff commits DATA blocks
using proposal timers and rotates leaders on terminal certificates.
A single DRAIN position is available only when recovery opens on a
full-quota chain. Cross-region replicas continuously fetch regional
blocks and finalized L1 control state. Normal completion retains the
terminal DATA prefix. An activated force instruction instead fixes an
external cutoff chosen by an honest multisignature committee. Old regional
blocks above that cutoff may remain locally committed but are excluded
from the external canonical history. Both paths use one bootstrap
handler. This draft specifies the model and protocols; security and
complexity analyses are deferred.
\end{abstract}

\section{Model and L1 control state}
\label{sec:model}
L1 is the public source of sequencing authority. It determines who may
vote and lead, how many DATA blocks a region may add, which regions are
eligible, and when an explicit forced rotation becomes effective.
Regional Sync HotStuff~\cite{sync} orders and commits blocks locally.
L1 force instructions determine which of those blocks remain in the
external canonical history.
Cross-region infrastructure distributes their contents, certificates, and
L1 proofs; it does not elect a leader by reporting a tip.

\subsection{What nodes read from L1}
\label{sec:l1}
A session $r$ is one region's production turn. Within it, $g$ counts
recoveries and $t$ counts leader terms. Nodes read the following L1 fields.

\begin{table}[htbp]\centering\small
\caption{L1 information used by the protocol.}
\label{tab:l1}
\begin{tabularx}{\linewidth}{@{}p{.27\linewidth}Y@{}}
\toprule
Field & Meaning\\
\midrule
Stake configuration & Regional members, public keys, voting weights,
and protocol parameters.\\
Leader schedule & $L(r,g,t)=(i_r,\ell_{r,g,t})$: the region and leader
assigned to a view.\\
Regional quota $TB_r$ & Maximum number of new DATA blocks in session
$r$. Local recovery does not reset it.\\
Region status $\chi$ & The live/off status of each region for future
selection.\\
Force flag & Whether this update triggers a forced rotation.\\
Target session $r$ & The session whose external history is to be
truncated, when the force flag is true.\\
Anchor height $h_*$ & The last DATA height retained from that session.
The update specifies a height, not a block hash.\\
Activation height $a$ & The update takes effect once L1 block $a$
is finalized.\\
\bottomrule
\end{tabularx}
\end{table}

The multisignature committee posts status updates containing the last
five fields in \cref{tab:l1}. When the force flag is false, only the status
map changes; target session and anchor height are unused.
A force update also marks its target region off for future selection.

Each update is announced at least the configured number of L1 blocks
before its activation height. It takes effect when that L1 block is
finalized. Updates have increasing activation heights and are applied
in L1 order. Nodes use the status map from the latest effective update.
They verify the update's authorization and L1 finality; finalized
records cannot be rewritten.

We write $\tagg{AUTH}_{r,g}$ for the configuration, leader schedule,
and quota used by generation $(r,g)$. These may be derived from preset
L1 rules, without a new transaction for every local rotation.
The initial configuration supplies the region order $\mathcal L_R$,
initial status map, and genesis. The first live region opens at genesis
with a synthetic regional receipt; if none is eligible, initialization
waits.

\begin{assumption}[L1 authority and control order]
\label{ass:authority}
L1 provides the configuration, leader schedule, and regional quota to
all correct replicas before the corresponding view begins, including
a recovery view. All correct replicas observe a common order of
finalized L1 updates and regional handoffs. L1 progresses independently
of regional BFT.
\end{assumption}

\subsection{Stake and service assumptions}
For region $i$, let $W_i=\sum_{v\in V_i}\omega_{i,v}$ with
$\omega_{i,v}>0$. Byzantine stake is at most $F_i<W_i/2$.
Every regional certificate in this draft requires distinct signers of
total weight strictly greater than $F_i$. Equivalently, for a fault-share
bound $f_i=F_i/W_i<1/2$, its weight exceeds $f_iW_i$.
Membership and weights are fixed throughout a session, including local
recoveries. Slashing or a passive status entry does not change its
denominator. A later session verifies old evidence against the old snapshot.
The bound applies to the cumulative compromised signing weight of each
snapshot for as long as its certificates are accepted; rotating corruption
or later compromise of more historical keys is outside this model.
Correct replicas retain their protocol state and committed blocks.
Signatures cannot be forged, and hash commitments are binding.

\begin{assumption}[Authenticated synchronous communication and global clock]
\label{ass:network}
\leavevmode\par
\begin{itemize}
    \item Messages use authenticated communication channels.
    \item Within region $i$, the actual delay of a correct-to-correct
    message is $\delta\le\Delta_i$; $\delta$ may differ across messages.
    \item All replicas share a global clock.
\end{itemize}
\end{assumption}
The delay bound includes delivery and processing of complete validation
dependencies. Each replica records its own entry and voting events on
the global clock.

\begin{assumption}[Cross-region communication backbone]
\label{ass:infra}
A communication backbone delivers the current region's committed block
stream, together with its regional finality certificates, to the next region.
\end{assumption}

\begin{assumption}[Honest multisignature committee]
\label{ass:cut}
The multisignature committee is honest.
\end{assumption}

\section{Views, blocks, and admission}
\label{sec:definitions}
A proposal QC records votes. A regional receipt records completed
commitment.

\subsection{Views and transitions}
A view is $(r,g,t)$: session $r$, recovery generation $g$, and leader
term $t$. Compare views by $r$ first, then $g$, then $t$.
A new regional session starts at $(r,0,0)$.

\begin{table}[htbp]\centering\small
\caption{The four ways to change view.}
\label{tab:views}
\begin{tabularx}{\linewidth}{@{}>{\raggedright\arraybackslash}p{.24\linewidth}>{\raggedright\arraybackslash}p{.23\linewidth}Y@{}}
\toprule
Event & Next view & Required event and timer rule\\
\midrule
Normal local rotation & $(r,g,t+1)$ & The term's terminal DATA has a
QC and is below the regional quota. Existing timers continue.\\
Local recovery & $(r,g+1,0)$ & Quit proof, recovery wait, and next
authority. Cancel unfinished timers in the old generation.\\
Regular regional handoff & $(r+1,0,0)$ & Terminal DATA has a regional
receipt, and the handoff fixes the successor.\\
Forced regional handoff & $(r+1,0,0)$ & L1 force has activated, and the
selected DATA anchor and handoff are verified. Old regional timers
are unchanged.\\
\bottomrule
\end{tabularx}
\end{table}
\FloatBarrier
A DRAIN QC causes no view change. A status update without an applicable
force instruction changes no view. Local recovery keeps the same
session quota.

\subsection{Blocks, anchors, and quotas}
\label{sec:positions}
The session anchor $A_r$ is the DATA block inherited from the previous
region; $A_0$ is genesis. It has a regional commitment receipt.
The generation opening $A_{r,g}$ is where local generation $g$ resumes.
Initially $A_{r,0}=A_r$. Recovery may choose a certified block that has
not yet committed. Once chosen, the opening stays fixed for that generation.

\paragraph{Two block positions.}
The global DATA height $h(B)$ counts application blocks.
The physical position $k_r(B)$ counts both DATA and DRAIN blocks within
session $r$. For this local count, set $k_r(A_r)=0$.
This convention does not change the inherited block's signed fields.

A new block has the form
\[
 B=(r,k,h,\mathit{type},x,\Hash(B_{\rm parent}),\rho),
 \qquad \mathit{type}\in\{\tagg{DATA},\tagg{DRAIN}\}.
\]
Here $x$ is the application batch and $\rho$ is a finalized L1 reference.
Every new block has $k=k_r(B_{\rm parent})+1$.
Its type determines the DATA height:
\begin{itemize}
\item DATA carries an application batch and has $h=h(B_{\rm parent})+1$.
\item DRAIN carries no application batch and has $h=h(B_{\rm parent})$.
\end{itemize}

\paragraph{The session quota.}
The session's terminal DATA height is
\[
 H_r=h(A_r)+TB_r.
\]
Write $T_r$ for the DATA block at $H_r$.
Count all DATA blocks in the proposed ancestry, including uncommitted
ones, when checking the quota. DRAIN, retransmission, and ANCHOR
re-certification do not consume DATA positions.

\paragraph{The leader-term quota.}
If the opening is below $H_r$, term $t$ can propose
\[
 m_{r,g,t}=\min\{K,H_r-h(A_{r,g})-tK\}>0
\]
new DATA blocks. Their positions $j=0,\ldots,m_{r,g,t}-1$ have
\[
 k=k_r(A_{r,g})+tK+j+1,\qquad
 h=h(A_{r,g})+tK+j+1.
\]
The last term may contain fewer than $K$ blocks.
No next term is admitted after its terminal DATA reaches $H_r$.

For example, start at DATA 100 with $TB_r=10$ and $K=4$.
The three terms produce DATA 101--104, 105--108, and 109--110.
If recovery instead opens on DATA 107, its first term has three remaining
positions: 108--110. Recovery does not reset the terminal height 110.

\paragraph{Restricted DRAIN.}
\label{sec:drain}
Recovery can inherit the terminal DATA before it commits.
Re-certifying this opening starts no timer. A fresh DRAIN proposal
provides the timer that can commit the opening and its ancestors.

DRAIN is allowed only under these conditions:
\begin{enumerate}[label=(\roman*)]
\item The generation is a recovery generation: $g>0$.
\item Its opening already has full DATA height: $h(A_{r,g})=H_r$.
The opening may be an older DRAIN.
\item Its only new position is $(t,j)=(0,0)$.
That DRAIN extends the opening using its current ANCHOR QC.
\end{enumerate}
The generation has no second proposal or next term.
A generation that first reaches $H_r$ with DATA cannot append DRAIN.
It waits for commitment or recovers after Quit.
Each later interrupted recovery may need its own DRAIN; the total across
generations is not bounded. Admission does not depend on whether an
unobserved replica has already committed the terminal DATA.

Committing DRAIN commits its DATA ancestors. The next region still
starts from the retained DATA block. DRAIN descendants remain in the old
region's control history.

\paragraph{L1 references.}
The reference $\rho$ contains a finalized L1 height $\lambda(B)$,
canonical block hash, and registry proof.
These heights do not decrease along DATA ancestry.
The first DATA reference is at least the session's opening control height.
DRAIN keeps its terminal DATA ancestor's reference.

Regional validation uses the session's fixed configuration and leader
authority. Later L1 updates do not rewrite this local validity test.
External validation also checks the current control history.
An old L1 reference cannot bypass an activated external cutoff.

\subsection{Which blocks the external chain retains}
\label{sec:external}
The external wrapper combines the DATA prefixes selected from regional
sessions. It uses the common control order assumed in Section 1.

A handoff is \emph{sealed} when that order fixes three things:
the retained DATA anchor, the successor region, and the L1 selection
height. Later updates cannot change this boundary.
Sealing requires no additional L1 transaction per regular handoff.

The wrapper tracks three states:
\begin{itemize}
\item \textbf{Open:} regional output is externally provisional, even if
the regional BFT has committed it.
\item \textbf{Truncated:} force has fixed the external cutoff.
The wrapper may still be waiting for anchor data or a successor.
\item \textbf{Sealed:} the retained prefix and successor are fixed.
The retained prefix is externally irrevocable.
\end{itemize}
The current external session remains $r$ until its handoff is sealed.
It then advances to the successor.

\paragraph{Normal and forced endpoints.}
Let $\operatorname{End}(r)$ be the session's external endpoint.
A sealed regular handoff sets $\operatorname{End}(r)=H_r$ and retains
$T_r$. Activated force sets $\operatorname{End}(r)=h_*$, even while
bootstrap is pending.

To request force, the multisignature committee records the force flag,
target session $r$, anchor height $h_*$, and activation height $a$.
It selects the last available committed DATA block in the region's
correctly operating prefix. That block, its regional receipt, and its
ancestry must already be available outside the old region.
The record specifies the height; the receipt identifies the block hash.
The cut must satisfy
\[
 h(A_r)\le h_*\le H_r.
\]
Thus it preserves the incoming sealed prefix.
Announcement allows preparation. L1 activation installs the cutoff.

\paragraph{External acceptance.}
For a DATA block $B$ produced in a closed session $r$, check:
\begin{enumerate}[label=(\roman*)]
\item Its regional receipt is valid:
$\operatorname{VerifyRegional}(B,\FC_r(B))=1$.
\item Its height satisfies $h(A_r)<h(B)\le\operatorname{End}(r)$.
\item The control history assigns its height to $r$.
The block is the retained anchor or one of its DATA ancestors.
This chain must extend $A_r$.
\end{enumerate}
For an open session, regional checks allow provisional output only,
up to $H_r$.

\paragraph{What force changes.}
On activation, discard the old session's provisional suffix above $h_*$.
Rebuild provisional application state from the retained prefix once its
data are available. Keep discarded admissible requests pending.
Do this even if successor startup must wait.

The old region may still propose or commit higher blocks.
Force does not change its timers, locks, $TB_r$, or committed log.
No shutdown proof or proof that higher blocks are absent is required.
Late receipts for those blocks remain regional evidence; they cannot
restore the discarded external suffix.
A later live status does not reopen the truncated session.

\paragraph{Requests and progress.}
Correct-client requests eventually reach the serving regions and external
pending pool. They have a common admission order and finitely many
predecessors. Each request fits in a batch.
Correct leaders include the oldest admissible requests that fit and are
not already in their proposed ancestry.
They may use empty DATA when no such request is available.

Remove a request from the external pool only after it enters a sealed
retained prefix. Recheck requests from a discarded suffix for inclusion
in the successor. Regional commitment alone does not remove them.

Progress also needs fair schedules and enough uninterrupted service to
produce DATA and complete handoffs.
Correct leaders must receive repeated turns in actually entered views,
including recovery initializers and DATA-producing turns across sessions.
Unreachable future assignments do not meet this condition.
Exposed equivocating leaders are eventually excluded by future authorities.
These are additional conditions for progress.

\subsection{Certificates and signed messages}
\label{sec:certificates}
Proposal QCs, ANCHOR QCs, and regional receipts each require matching
signatures from distinct members with total weight $>F_{i_r}$.

\begin{table}[htbp]\centering\small
\caption{Votes, anchor approval, and commitment evidence.}
\label{tab:certificates}
\begin{tabularx}{\linewidth}{@{}>{\raggedright\arraybackslash}p{.25\linewidth}>{\raggedright\arraybackslash}p{.23\linewidth}Y@{}}
\toprule
Certificate & Signed message & Purpose\\
\midrule
Proposal QC $\QC_{r,g,t}(B)$ & VOTE for DATA or DRAIN & Authorize a
permitted child or normal term entry.\\
ANCHOR QC & ANCHOR-VOTE & Certify the unchanged generation opening
before the first fresh proposal.\\
Regional receipt $\FC_r(B)$ & FINAL for DATA & Prove regional
commitment. External acceptance still follows the selected boundary.\\
\bottomrule
\end{tabularx}
\end{table}

A correct replica casts at most one fresh vote per $(r,g,t,j)$ and
one ANCHOR vote per $(r,g)$.
DATA and DRAIN use the same fresh-vote rule but have distinct signed types.
A QC or an ANCHOR vote starts no commit timer.

\paragraph{Certificate rank and the base anchor.}
Within a session, rank certificates by $(g,t,k_r(B))$, in that order.
For equal ranks, use a fixed total tie-break on the certified content.
Signer subsets and signature encodings do not affect the tie-break.
An opening's identity includes its hash and positions.

The inherited anchor has a base justification $Q_r^0$.
It ranks below every ordinary QC in this session.
Recovery can use it if no ordinary QC exists.
An old region's higher-ranked QC cannot replace the selected incoming
anchor.

\paragraph{Signing fields.}
Every message binds the chain, voting configuration, session, region,
and message type. View-dependent messages also bind the L1 authority
and $(r,g,t)$. Additional fields are:
\begin{itemize}
\item \textbf{PROPOSE and fresh VOTE:} opening, position $j$, block type,
physical position, DATA height, and block hash.
\item \textbf{NEWVIEW and ANCHOR-VOTE:} the selected opening and its block
identity. NEWVIEW also binds the justification and entry proof.
\item \textbf{STATUS:} the sender's fixed entry-lock certificate and
entry proof. The next opening need not be known.
\item \textbf{BLAME:} the admitted view, without an opening, block hash,
or local progress count. All complaints for that view sign the same
statement, even before NEWVIEW.
\end{itemize}

\paragraph{Regional receipts.}
After committing a DATA block $B$ produced in session $r$, a correct
replica sends
\[
 \langle\tagg{FINAL},\mathit{cfg}_r,r,h(B),\Hash(B)\rangle_v
\]
to peers and the backbone. This applies whether $B$ commits on its own
timer or through a descendant. Matching endorsements form $\FC_r(B)$.
The receipt includes $B$ and the needed ancestry, or resolves them
from a validated cache.

FINAL binds the chain and these session-stable fields.
It does not bind the generation or term in which the signer commits.
Confirmations of the same block can therefore combine across recoveries.
The block hash preserves its original signed context.

If DRAIN commits $T_r$, export the receipt for $T_r$, not for DRAIN.
The incoming anchor keeps its historical receipt.
A replica may adopt a valid receipt for its own session to commit that
DATA prefix locally. Other sessions' receipts provide history or anchor
evidence; they do not replace the current engine's locks or log.

\paragraph{Parent selection.}
\Cref{tab:parent} gives the parent and certificate for each fresh position
in an admitted generation.

\begin{table}[htbp]\centering\small
\caption{Parent selection within an admitted generation.}
\label{tab:parent}
\begin{tabularx}{\linewidth}{@{}p{.32\linewidth}Y Y@{}}\toprule
Position & Parent & Parent certificate\\\midrule
First DATA, $t=j=0$ & $A_{r,g}$ & ANCHOR QC in $(r,g,0)$\\
First DATA, $t>0,j=0$ & Previous term's terminal DATA & Its proposal QC\\
Later DATA, $j>0$ & Position $j-1$ of this term & Its proposal QC\\
Only DRAIN, $g>0,t=j=0$ & Full-quota opening $A_{r,g}$ & ANCHOR QC in $(r,g,0)$\\
\bottomrule\end{tabularx}
\end{table}
\FloatBarrier

\subsection{Entry, message processing, and regional state}
\label{sec:admission}
\paragraph{What admits a view.}
The entry evidence depends on the transition:
\begin{itemize}
\item \textbf{New region, $g=0,t=0$:} Ready and the boundary proof admit
the view. NEWVIEW must select that boundary's exact anchor.
\item \textbf{Local recovery, $g>0,t=0$:} the quit-entry proof, recovery
wait, and next authority admit the view.
NEWVIEW uses an older-generation QC from the same session, or $Q_r^0$.
\item \textbf{Normal term, $t>0$:} all preceding term-terminal QCs in
the generation must be present. Each certifies DATA below $H_r$.
\end{itemize}
Term zero is admitted before NEWVIEW arrives.
Its progress clock therefore runs even if the initial leader stays silent.
An ANCHOR QC is required before fresh proposals begin.
Genesis is the sole synthetic initial anchor.

\paragraph{How to process a proposal.}
\begin{enumerate}[label=(\roman*),itemsep=2pt]
\item Check the available signed material for valid quit evidence.
Do this before rejecting an outer message or discarding a duplicate.
Compare incoming material with retained history.
\item Validate signatures, authority, entry evidence, type, quota,
both positions, parent QCs, hash links, and application validity.
Required ancestors include their signed proposals and block data.
Correct senders supply complete bundles; cached references must resolve
within the complete-message bound.
\item In the current generation, invoke Vote for each admitted DATA or
DRAIN proposal not yet voted on. This includes earlier-term proposals,
attached history, and blocks already committed as ancestors.
Process ancestors before children and before terminal-QC entry.
\item Deduplicate by $(r,g,t,j,\Hash(B))$ after inspecting evidence.
Retain old-generation proposals without new votes or timers.
Buffer future-generation proposals until entry.
\end{enumerate}
At equal timestamps, process received messages before timer callbacks.

\paragraph{Regional state.}
Each regional engine has its own $(r,g,t)$ and one of three states:
\begin{itemize}
\item \textbf{Active:} an admitted, unfrozen generation may propose,
vote, and make normal term transitions while the session is incomplete.
\item \textbf{Frozen:} stop new proposals and votes, but still process
evidence, old votes, QCs, regional receipts, and L1 updates.
\item \textbf{Complete:} a verified receipt for terminal DATA ends
regional production, even if the external handoff is still pending.
\end{itemize}
Local terminal commitment alone is not enough to stop participating.
Other members may still need recovery votes before a terminal receipt forms.
The external wrapper's open, truncated, and sealed states are separate.
Its cursor may advance while an old regional engine continues.
Route evidence by session so that an old session cannot freeze its
successor's engine.

\section{Protocols}
\label{sec:protocols}
Protocols 1 and 2 patch v3.1's steady-state and recovery boxes. Red text
marks substantive changes from v3.1, including those retained from earlier RR drafts;
the session index $r$ and regional delay subscript are implicit throughout.
Protocol 3 is new relative to v3.1 and is shown in red. Explanations and service assumptions are
kept outside the execution steps.

\begin{protocol}{alg:steady}{Propose, vote, commit, and regular termination}
\begin{enumerate}[label=\textbf{\arabic*)},itemsep=5pt]
\item \textbf{Propose.} The authorized leader $\ell_{r,g,t}$ obtains the
parent and its QC from \cref{tab:parent}. \RR{If this is a DATA generation,
use the next position $j<m_{r,g,t}$. If it is a drain generation, use only
its one DRAIN position. A regionally complete session proposes nothing.}
Broadcast the complete proposal
\[
 P=\langle\tagg{PROPOSE},r,g,t,j,B,Q_{\rm parent}\rangle_{\ell_{r,g,t}}
\]
and process it locally under Vote. The parent need not be committed.

\item \textbf{Vote.} Process quit evidence first; return if it freezes
the current generation. Check completeness and validity. Request missing
data without voting; reject invalid proposals. If already voted for $P$,
stop. Otherwise forward $P$, broadcast
\RR{$\langle\tagg{VOTE},r,g,t,j,\Hash(B)\rangle_v$}, and start its
timer at the local acceptance time $\tau_v(P)$, expiring at
\RR{$\tau_v(P)+2\Delta_{i_r}$}. Process unseen earlier-term proposals
and attached history by the same rule. A duplicate never restarts a timer.

\item \textbf{Commit and export.} On timer expiry, commit $B$ and its
uncommitted ancestors. Normal term rotation preserves timers; Quit cancels
unfinished generation timers. A QC alone starts no timer.
\RR{Send FINAL endorsements for newly committed DATA produced in this
session, including the terminal DATA ancestor of a committed DRAIN.
Relay regional receipts to peers and the cross-region infrastructure;
external acceptance is determined by Protocol 3.}

\item \textbf{Terminal QC.} On a valid term-terminal proposal QC, process
its unseen history under Vote. \RR{If it certifies DATA strictly below
$H_r$}, and the generation remains unfrozen, send its complete bundle
once to $\ell_{r,g,t+1}$ and enter $(r,g,t+1)$ if not already there or
beyond it. Record first entry $e_v(r,g,t+1)$ and disable completed-term
checks. Process terminal-QC chains in order; duplicates never reset entry.
\RR{A QC for DATA at $H_r$, or for DRAIN, admits no next term or fresh
position. Continue pending timers and evidence processing.}

\item \RR{\textbf{Complete the regional engine.} On a valid $\FC_r(T_r)$,
adopt its DATA prefix, relay the terminal bundle, and mark the engine for
session $r$ complete. Stop its new proposals, votes, progress checks,
unfinished timers, and recovery. Continue serving data and running
Protocol 3. Force instructions do not invoke this step or cancel
regional timers; Protocol 3 applies their external cutoff.}
\end{enumerate}
\end{protocol}

Acceptance occurs only after complete validation and the first vote.
The term-entry clock starts at admission, before any first proposal by
that term's leader. Reaching $H_r$ in a proposal or QC is not a cross-region
entry event. With no Quit, the terminal DATA's own timer is sufficient;
no subsequent block is required for that commitment.

\paragraph{Complete quit proofs.}
For an admitted $(r,g)$, a proof is either: (a) matching
$\langle\tagg{BLAME},r,g,t\rangle$ signatures of total weight
$>F_{i_r}$ for an admitted term, including admitted term zero before
NEWVIEW; (b) two different admissible fresh
proposals signed by its leader at the same $(r,g,t,j)$; or (c) two
admissible NEWVIEWs signed by $\ell_{r,g,0}$ selecting different openings.
The proposals include their admission and quota witnesses. An excess
DATA block or unauthorized DRAIN is rejected, not treated as a second
admitted position. Different openings constitute equivocation even if
comparable; different certificates for the same opening do not.
Evidence from any admitted term can freeze its generation. Old-session
or old-generation evidence is archival and cannot freeze a current
successor. Future proofs wait for the corresponding authenticated entry.

\begin{protocol}{alg:recovery}{Blame, quit, and recovery with restricted drain}
\begin{enumerate}[label=\textbf{\arabic*)},itemsep=5pt]
\item \textbf{Blame.} At first term entry, schedule the first check at
\RR{$e_v(r,g,t)+6\Delta_{i_r}$}, even before NEWVIEW or a proposal arrives.
\RR{Once the opening is authenticated, use $q=m_{r,g,t}$ for a DATA term
and $q=1$ for a drain generation.} For $p=1,\ldots,q$, check at
\[
 e_v(r,g,t)+(2p+4)\Delta_{i_r}.
\]
If the term remains current and unfrozen and fewer than $p$ distinct
fresh positions have triggered local votes, broadcast
$\langle\tagg{BLAME},r,g,t\rangle_v$ once.
\RR{The one DRAIN vote counts; ANCHOR votes do not. Check an already-due
deadline when its opening becomes known. Do not require positions beyond $q$.}

\item \textbf{Quit.} On the first complete proof $E$ for current $(r,g)$,
freeze the generation: stop its progress checks, new proposals, votes,
and normal transitions; cancel all unfinished proposal timers.
Broadcast the proof and record quit time $a_v$ once. Retain commitments
and process old votes and QCs. Submit equivocation to L1 without awaiting
its finalization. \RR{Continue receiving regional receipts; the external watcher
processes L1 controls independently.}

\item \textbf{Status and entry.} \RR{Unless the regional engine has
completed its quota}, enter at
\[
 \RR{e_v(r,g+1,0)=a_v+\Delta_{i_r}.}
\]
\RR{The timely-authority assumption makes $\tagg{AUTH}_{r,g+1}$
available before this entry.} Read that authority. Fix the highest-ranked retained lock or valid
known QC from generations through $g$ as entry lock $Q_v$
\RR{(including the retained base justification $Q_r^0$)}.
Enter $(r,g+1,0)$, mark it unfrozen, start its progress clock, and send
$\langle\tagg{STATUS},r,g+1,0,Q_v,E\rangle_v$ with complete dependencies
to $\ell_{r,g+1,0}$. Process buffered proofs immediately. Later QCs update
retained evidence, not this fixed lock. \RR{Keep the original $A_r$,
$H_r$, and voting snapshot.}

\item \textbf{New-view.} The leader waits $2\Delta_{i_r}$ from its own
entry. If still current and unfrozen, choose the highest valid known QC
$Q^*$ from generations through $g$ \RR{or the base justification if no
ordinary QC is known}, and broadcast
\[
 \langle\tagg{NEWVIEW},r,g+1,0,A,Q^*,E\rangle_{\ell_{r,g+1,0}},
\]
where $Q^*$ certifies $A$, with complete dependencies. Choose one opening
for this generation. \RR{An opening below $H_r$ resumes DATA; an opening
at $H_r$, including a DRAIN descendant, admits exactly one fresh DRAIN.}

\item \textbf{First vote.} \RR{This handler applies to any admitted
term-zero view $(r,u,0)$, including bootstrap $u=0$.} Process quit evidence,
then validate NEWVIEW and its dependencies. If current and unfrozen,
no ANCHOR vote was cast, and $\rank(Q^*)\ge\rank(Q_v)$, forward NEWVIEW
and broadcast
$\langle\tagg{ANCHOR\mbox{-}VOTE},r,u,0,\Hash(A)\rangle_v$.
Start no timer and consume no fresh position. On an ANCHOR QC
\RR{of weight $>F_{i_r}$}, the leader resumes Protocol 1 at $j=0$.
\end{enumerate}
\end{protocol}

The entry-lock test governs ANCHOR voting only. A subsequent fresh
proposal backed by a valid ANCHOR QC is processed under Protocol 1 even
at a replica that declined the ANCHOR vote. A locally committed terminal
block does not justify stopping this recovery before a terminal regional
receipt is available. An external force cutoff is not a local quit proof
and does not alter this recovery rule.

\subsection{Boundary witnesses and the bootstrap event}
\label{sec:boundary}
For an effective map $\chi$, let $\tagg{NextLive}(i,\chi)$ be the first
other live region after $i$ in $\mathcal L_R$, searching at most one cycle.
If none exists, its value is $\bot$. A boundary witness has the form
\[
 \mathcal W=(\mathit{mode},r,r+1,i',A,\FC(A),\lambda_s,\Xi),
 \qquad \mathit{mode}\in\{\tagg{REGULAR},\tagg{FORCE}\},
\]
where $\Xi$ authenticates the control order, effective map at selection
height $\lambda_s$, and any activated L1 force entry. It also
binds the successor authority and configuration. The anchor's hash is
fixed by the block and finality receipt inside $\mathcal W$; the
L1 update itself specifies only the anchor height. The accepted witness seals that successor,
retained DATA prefix, and selection height. Certificate encodings do not
change boundary identity. A force witness requires no evidence that the
old region has stopped or lacks higher regional commitments.
Every external verifier can validate this witness and install its sealed
prefix; membership in the successor region is needed only to bootstrap
that region's engine.

For a regular boundary, $A=T_r$, its ancestry contains exactly $TB_r$
DATA blocks since $A_r$, and selection starts at the terminal block's
signed L1 height. For a forced boundary, $A$ is the regionally committed DATA block
at $h_*$ named by the activated entry, with
$h(A_r)\le h_*\le H_r$. If $h_*=h(A_r)$, reuse the incoming anchor and
its historical receipt. No new receipt from the failed region is required
for a prefix it inherited unchanged.
For force, the selection height $\lambda_s$ is at or after both the
activation height $a$ and the anchor's L1 reference $\lambda(A)$.
Apply all updates effective by $\lambda_s$ before choosing the successor.
For genesis, the initial L1 reference supplies $\lambda(A_0)$.

At an unsealed boundary, process due controls in their authenticated
order and apply an applicable force before regular quota completion at
the same control height. A sealed boundary is immutable. A force targeting
an already departed session cannot replace its successor; its status map
still applies in the ordered control history. A later live update does
not erase the installed cutoff or reopen that session's external
authorization.
If $i'=\bot$, retain the anchor and advance the authenticated control
cursor; the first admissible later prefix with an eligible successor seals
the waiting boundary. A newer local L1 head cannot reselect it afterward.

\begin{definition}[Local cross-region bootstrap event]
\label{def:bootstrap}
For a replica $v\in V_{i'}$, $\tagg{Ready}_v(\mathcal W)$ occurs on its
first complete validation of an admissible sealed witness naming $i'$,
the anchor and missing application history, and the matching finalized
$\tagg{AUTH}_{r+1,0}$. REGULAR requires terminal regional finality and the regional
quota; FORCE additionally requires finalized L1 through its activation
height and the corresponding external cutoff in the ordered control history. The event is suppressed if this replica
already installed that successor or has authenticated a later session.
An incomplete boundary stays pending until its data and authority arrive.
Prospective members process boundary bundles before regional entry;
other new-session protocol messages wait for that authenticated entry.
\end{definition}
This is the precise event at which a node recognizes that its region may
start the next session. Merely receiving a leading block at height $H_r$,
its QC, an announced force entry, or a newer local status map is insufficient.
Ready starts the local entry clock; a remote proposal time, private QC
assembly, or L1 installation time does not.
On first complete validation, the recipient immediately forwards the
bundle to all members of the selected region. The bundle supplies all
data and proofs needed by its receivers to validate the boundary and
reconstruct the anchor state; cached references must resolve within the
regional complete-message bound. Receivers validate and enter without
waiting for another backbone delivery or a leader message.

\paragraph{Terminal recovery example.}
Let $h(A_r)=100$, $TB_r=3$, and $K=2$. The first term proposes DATA
101--102 and its terminal QC admits the next term, which has only DATA
103. Its QC starts no next term. If its own timer survives, the receipt
for DATA 103 enables the regular handler. If Quit cancels the timers and
the next generation opens on certified DATA 103, its sole fresh proposal
is one DRAIN. That DRAIN's timer commits DATA 103 as an ancestor and
produces the same terminal DATA receipt. Another Quit may require a new
generation and one further DRAIN. In every case, the next region starts
with DATA 104, whose parent is DATA 103.

\paragraph{Forced bootstrap example.}
Suppose a session opened at DATA 100 with $TB_r=10$. A multisignature
entry specifies cutoff 106 and force activation at L1 height 1200.
The old region may locally commit DATA 107 and 108 during the notice
period or afterward. Once L1 through 1200 is finalized, external nodes
retain this session only through DATA 106. They discard any provisional
old-session suffix above 106 and recheck its requests for inclusion.
After verifying DATA 106 and its regional receipt, the selected successor
opens $(r+1,0,0)$ and proposes its own DATA 107 with parent DATA 106.
The old regional log is unchanged. Its receipts for 107 and 108 do not
restore those blocks to the external canonical chain.

\begin{protocol}{alg:bootstrap}{Watch the stream and bootstrap the next region}
\color{rrRed}
\begin{enumerate}[label=\textbf{\arabic*)},itemsep=5pt]
\item \textbf{Watch and fetch.} Prospective successors continuously
receive the current session's committed stream and regional receipts
through the backbone. Cache earlier proposals or QCs for preparation,
fetch missing ranges, and independently poll finalized L1.
Validate and retain the handoff data and proofs. Apply ordered controls
to external output: discard suffixes above installed cutoffs regardless
of old L1 references, and treat open-session output as provisional.
Any external verifier receiving a valid sealed witness installs its
retained prefix, endpoint, and successor and updates the pending pool.

\item \textbf{Handle regular completion.} Verify the terminal DATA
receipt and exactly $TB_r$ DATA positions. Replay boundary controls from
its signed L1 reference. If no force has taken precedence, choose
$i'=\tagg{NextLive}(i_r,\chi)$ using the status map effective at the
recorded selection height $\lambda_s$, and verify the regular witness
through $H_r$. If none is eligible, retain the candidate anchor and
watch later controls. A regional receipt cannot override a force cutoff.

\item \textbf{Handle activated force.} Once L1 block $a$ is finalized,
verify an authorized update with force flag true, target session $r$,
and anchor height $h_*$. Reject a stale target. For the current external
session, install $\operatorname{End}(r)=h_*$ once and discard its higher
provisional suffix. Fetch the anchor, historical regional receipt, and
ancestry; restore the retained application state and keep discarded
admissible requests pending. Regional timers, locks, and commitments
are unchanged. Choose the successor using all controls effective at
$\lambda_s$, which must be at or after both $a$ and $\lambda(B_{h_*})$.
If data or a successor is missing, keep the cutoff and wait.

\item \textbf{Enter on Ready.} On $\tagg{Ready}_v(\mathcal W)$, immediately
relay the complete validation bundle to all members of the selected
region. Receivers validate and enter without waiting for another backbone
delivery or leader message. Install session $r+1$, exact anchor $A_{r+1}=A$,
configuration, quota, and opening height $\lambda_s$. Apply the external
update in step 1. Set the retained and entry lock $Q_v=Q_{r+1}^0$.
Enter $(r+1,0,0)$ unfrozen, record $e_v(r+1,0,0)$ once, and start its
progress clock. Send STATUS with the base justification and $\mathcal W$
to $\ell_{r+1,0,0}$; process buffered current-generation evidence.
Duplicates never reset entry, quota, or locks.

\item \textbf{Announce and certify the opening.} After $2\Delta_{i'}$
from its entry, the initial leader, if still current and unfrozen,
broadcasts NEWVIEW with exactly $A_{r+1}$, $Q^*=Q_{r+1}^0$, and
$\mathcal W$. Members validate the anchor, authority, and entry proof,
then forward and ANCHOR-vote under Protocol 2's first-vote rule.
ANCHOR voting starts no timer.

\item \textbf{Start the region.} On an ANCHOR QC of weight $>F_{i'}$,
the initial leader invokes Protocol 1 at DATA height $h(A_{r+1})+1$
and local physical position $1$, with quota $TB_{r+1}$.
Local failures use Protocol 2 within the same session.
Continue watching L1 and streaming blocks and regional receipts.
\end{enumerate}
\end{protocol}

Both handlers rerun on stream, authority, or control updates, including
while the engine is frozen, and are idempotent. Before entry or production,
apply all controls due for the successor. Genesis uses a synthetic
boundary with the same exchange.

\paragraph{Bootstrap timing.}
Let $e$ be the first correct Ready time. While the boundary remains current,
the complete-bundle relay gives entry everywhere by $e+\Delta_{i'}$.
With a correct initializer and no interruption, NEWVIEW is sent by
$e+3\Delta_{i'}$, its ANCHOR QC reaches the leader by $e+5\Delta_{i'}$,
and first DATA reaches all correct members by $e+6\Delta_{i'}$, meeting
the existing deadline. Backbone transfer before Ready has no fixed bound here.

\begingroup\small
\begin{thebibliography}{9}
\bibitem{sync}
I. Abraham, D. Malkhi, K. Nayak, L. Ren, and M. Yin.
\emph{Sync HotStuff: Simple and Practical Synchronous State Machine Replication}.
IEEE Symposium on Security and Privacy, 2020, pp. 106--118.
\href{https://conferences.computer.org/sp/pdfs/sp/2020/349700a106.pdf}{Published paper}.
\end{thebibliography}
\endgroup
\end{document}
